\documentclass[11pt]{article}

\usepackage[margin=1in]{geometry}
\usepackage{amsmath,amssymb,amsthm}
\usepackage{mathtools}
\usepackage{booktabs}
\usepackage{hyperref}
\usepackage{xcolor}
\usepackage{enumitem}
\usepackage{titlesec}
\usepackage{fancyhdr}

\hypersetup{colorlinks=true, linkcolor=blue!50!black, citecolor=blue!50!black,
            urlcolor=blue!50!black}

\titleformat{\section}{\large\bfseries}{\thesection}{0.6em}{}
\titleformat{\subsection}{\normalsize\bfseries}{\thesubsection}{0.6em}{}

\pagestyle{fancy}
\fancyhf{}
\rhead{\small Geographic Arbitrage in Physical Commodity Markets}
\lhead{\small \thepage}

\theoremstyle{definition}
\newtheorem{definition}{Definition}
\newtheorem{assumption}{Assumption}
\theoremstyle{plain}
\newtheorem{proposition}{Proposition}
\newtheorem{theorem}{Theorem}
\newtheorem{lemma}{Lemma}
\newtheorem{corollary}{Corollary}

\newcommand{\E}{\mathbb{E}}
\newcommand{\R}{\mathbb{R}}
\newcommand{\Var}{\operatorname{Var}}
\newcommand{\Cov}{\operatorname{Cov}}
\newcommand{\inprob}{\xrightarrow{p}}
\newcommand{\indist}{\xrightarrow{d}}

\title{\textbf{The Mathematics of Geographic Arbitrage\\ in Physical Commodity Markets}\\[4pt]
\large A Rigorous Treatment with a Synthetic-Data Laboratory}
\author{\textit{geo-arb} project documentation}
\date{\today}

\begin{document}
\maketitle

\begin{abstract}
\noindent
Physical commodity merchants such as the large agricultural trading houses exist,
in large part, to transform commodities across \emph{space}, \emph{time}, and
\emph{form}. This document gives a self-contained, rigorous account of the
\emph{spatial} (geographic) transformation: buying a commodity at one location
and selling it at another. We formalise the delivered-cost identity, the
spatial no-arbitrage band, and the Takayama--Judge spatial price equilibrium as
a linear program whose dual variables are equilibrium location prices. We then
specify a fully controlled stochastic data-generating process in which regional
prices are \emph{cointegrated} through a common world factor, freight and basis
are mean-reverting, and FX follows a random walk. Within this process we derive
the distribution of the arbitrage margin and the conditions under which a
window opens. The companion \texttt{geoarb} Python package implements every
object defined here, so each proposition can be checked numerically against
simulation. \textbf{All data are synthetic; nothing here is investment advice.}
\end{abstract}

\tableofcontents
\bigskip

%======================================================================
\section{Setup: the spatial market}
\label{sec:setup}

We consider a single homogeneous commodity (say soybeans) traded across a
finite set of locations partitioned into \emph{origins} (export ports) and
\emph{destinations} (import ports).

\begin{definition}[Spatial network]
Let $\mathcal{O}=\{1,\dots,m\}$ index origins and $\mathcal{D}=\{1,\dots,n\}$
index destinations. A \emph{lane} is an ordered pair $(i,j)\in\mathcal{O}\times
\mathcal{D}$. Each lane has a great-circle distance $\delta_{ij}>0$ (nautical
miles) and a voyage time $\tau_{ij}=\delta_{ij}/(24\,v)$ days for a vessel of
speed $v$ knots.
\end{definition}

At each trading date $t\in\{0,1,\dots,T\}$ we observe:
\begin{itemize}[topsep=2pt,itemsep=1pt]
  \item origin local cash prices $P^{o}_{i,t}$ and destination cash prices
        $P^{d}_{j,t}$ in USD per tonne;
  \item per-tonne freight $f_{ij,t}$ on each lane;
  \item port handling costs $h_i^o,\;h_j^d$;
  \item an import tariff $\theta_{ij}\ge 0$;
  \item an FX rate $e_{i,t}$ (USD per unit of origin $i$'s local currency);
  \item a benchmark futures price $H_t$ on a reference exchange.
\end{itemize}

%======================================================================
\section{Delivered cost and the arbitrage margin}
\label{sec:delivered}

The economic primitive is the \emph{delivered cost}: the all-in per-tonne cost,
expressed at the destination, of sourcing the commodity at origin $i$ and
landing it at destination $j$ at the time the cargo arrives.

\begin{definition}[Delivered cost]
\label{def:delivered}
With financing rate $r$ (annualised, $360$-day convention), the delivered cost
of lane $(i,j)$ initiated at date $t$ is
\begin{equation}
C_{ij,t}
=\underbrace{\frac{P^{o}_{i,t}}{e_{i,t}}+h^o_i}_{\text{FOB cost}}
+\;f_{ij,t}+\theta_{ij}+h^d_j
+\;\underbrace{\Big(\tfrac{P^o_{i,t}}{e_{i,t}}+h^o_i\Big)\, r\,\tfrac{\tau_{ij}}{360}}_{\text{financing of cargo in transit}}.
\label{eq:delivered}
\end{equation}
\end{definition}

\begin{definition}[Arbitrage margin / window]
The (gross) \emph{arbitrage margin} of lane $(i,j)$ at date $t$ is
\begin{equation}
M_{ij,t}\;=\;P^{d}_{j,t}\;-\;C_{ij,t}.
\label{eq:margin}
\end{equation}
The lane's \emph{arbitrage window} is \emph{open} at $t$ if $M_{ij,t}>0$ and
\emph{closed} otherwise.
\end{definition}

The margin (\ref{eq:margin}) is the per-tonne profit a trader books by buying
FOB at $i$ and selling at $j$ today, ignoring price moves during transit;
Section~\ref{sec:risk} restores that risk.

%======================================================================
\section{The spatial no-arbitrage band}
\label{sec:noarb}

\begin{assumption}[Frictional law of one price]
\label{ass:lop}
Capital is mobile and the number of potential shippers is large, so any lane
with a strictly positive, riskless margin is exploited until the margin is
competed away.
\end{assumption}

\begin{proposition}[Spatial no-arbitrage band]
\label{prop:band}
Under Assumption~\ref{ass:lop}, in a frictionless competitive equilibrium with
no binding capacity constraints,
\begin{equation}
P^{d}_{j,t}-P^{o}_{i,t}/e_{i,t}\;\le\; h^o_i+f_{ij,t}+\theta_{ij}+h^d_j
+\Big(\tfrac{P^o_{i,t}}{e_{i,t}}+h^o_i\Big) r\tfrac{\tau_{ij}}{360}
\qquad\forall (i,j),
\label{eq:band}
\end{equation}
equivalently $M_{ij,t}\le 0$ for every lane, with equality on every lane that
carries strictly positive flow.
\end{proposition}

\begin{proof}
Suppose $M_{ij,t}>0$ for some lane. A trader shipping one extra tonne earns
$M_{ij,t}>0$ at no risk, contradicting equilibrium under
Assumption~\ref{ass:lop} unless the action is bounded. Competition drives
$P^d_{j,t}$ down and $P^o_{i,t}$ up until $M_{ij,t}=0$. Lanes that remain
strictly unprofitable ($M_{ij,t}<0$) carry no flow. Hence at equilibrium every
active lane satisfies $M_{ij,t}=0$ and every lane satisfies
$M_{ij,t}\le 0$, which is (\ref{eq:band}). \qed
\end{proof}

Equation~(\ref{eq:band}) is the \emph{spatial} version of the Law of One Price:
prices at two locations may differ, but only up to the full cost of moving
product between them. The width of the band is precisely the right-hand side of
(\ref{eq:band}) --- the transfer cost --- and \emph{not} zero. This is why
persistent price gaps between continents are not, by themselves, evidence of
inefficiency.

%======================================================================
\section{Spatial price equilibrium as a linear program}
\label{sec:spe}

When capacity binds and many origins compete for many destinations, the
equilibrium flow pattern solves the classical Takayama--Judge transportation
problem.

\begin{definition}[Spatial price equilibrium, LP form]
Given supplies $s_i$, demands $d_j$, and delivered costs $C_{ij}$ at a fixed
date, the competitive equilibrium flows $x_{ij}\ge 0$ solve
\begin{equation}
\begin{aligned}
\max_{x\ge 0}\quad & \sum_{i\in\mathcal O}\sum_{j\in\mathcal D}\bigl(P^d_j-C_{ij}\bigr)x_{ij}\\
\text{s.t.}\quad & \sum_{j} x_{ij}\le s_i &&(\lambda_i)\\
                 & \sum_{i} x_{ij}\le d_j &&(\mu_j)
\end{aligned}
\label{eq:lp}
\end{equation}
where $\lambda_i,\mu_j\ge 0$ are the dual multipliers.
\end{definition}

\begin{theorem}[Duals are equilibrium prices]
\label{thm:dual}
Let $x^\star,(\lambda^\star,\mu^\star)$ be primal--dual optimal for
(\ref{eq:lp}). Then complementary slackness gives, for every lane,
\begin{equation}
P^d_j - C_{ij}\;\le\;\lambda^\star_i+\mu^\star_j,
\qquad\text{with equality whenever } x^\star_{ij}>0.
\label{eq:cs}
\end{equation}
The shadow price $\lambda^\star_i$ is the marginal value of an extra tonne of
supply at origin $i$, and $\mu^\star_j$ the marginal value of an extra tonne of
demand at destination $j$. Active lanes earn exactly the spatial price
difference implied by the duals; inactive lanes are weakly unprofitable.
\end{theorem}

\begin{proof}
The Lagrangian of (\ref{eq:lp}) is
$\mathcal L=\sum_{ij}(P^d_j-C_{ij})x_{ij}
-\sum_i\lambda_i(\sum_j x_{ij}-s_i)-\sum_j\mu_j(\sum_i x_{ij}-d_j)$.
Stationarity in $x_{ij}\ge 0$ (KKT) requires
$P^d_j-C_{ij}-\lambda_i-\mu_j\le 0$, with equality when $x_{ij}>0$, which is
(\ref{eq:cs}). The envelope theorem gives
$\partial(\text{optimal value})/\partial s_i=\lambda^\star_i$ and
$\partial/\partial d_j=\mu^\star_j$, the stated marginal valuations. \qed
\end{proof}

Theorem~\ref{thm:dual} makes Proposition~\ref{prop:band} precise in the
capacity-constrained case: the no-arbitrage band holds \emph{lane by lane} with
the duals replacing raw prices, and only the cheapest-to-serve origin--
destination pairs carry flow. This is the mechanism by which competition
``closes'' arbitrage windows even when a naive scan of prices shows many
positive margins.

%======================================================================
\section{A controlled stochastic data-generating process}
\label{sec:dgp}

To study \emph{when} windows open we need a model of how prices, freight and FX
co-move. We use a common-trend (factor) model that produces cointegrated
regional prices, the empirically dominant feature of integrated commodity
markets.

\subsection{Prices: a common stochastic trend}

\begin{definition}[World factor and location prices]
Let the world factor follow a random walk with drift,
\begin{equation}
F_t = F_{t-1}+\nu+\eta_t,\qquad \eta_t\sim\mathcal N(0,\sigma_F^2)\ \text{iid}.
\end{equation}
Each location price is the common trend plus a constant intercept $a_k$ plus a
stationary idiosyncratic spread $u_{k,t}$:
\begin{equation}
P_{k,t}=F_t+a_k+u_{k,t},\qquad
u_{k,t}=\rho\,u_{k,t-1}+\varepsilon_{k,t},\quad |\rho|<1,
\label{eq:price}
\end{equation}
with $\varepsilon_{k,t}\sim\mathcal N(0,\sigma_u^2)$ iid across $k,t$.
\end{definition}

\begin{proposition}[Regional prices are cointegrated]
\label{prop:coint}
For any two locations $k,\ell$, the spread
$Z_{t}:=P_{k,t}-P_{\ell,t}$ is (covariance-)stationary; equivalently $P_{k}$ and
$P_{\ell}$ are cointegrated $CI(1,1)$ with cointegrating vector $(1,-1)$.
\end{proposition}

\begin{proof}
From (\ref{eq:price}), $Z_t=(a_k-a_\ell)+(u_{k,t}-u_{\ell,t})$. The common
trend $F_t$ cancels. The difference of two independent stationary AR(1)
processes is stationary, so $Z_t$ is stationary with mean $a_k-a_\ell$. Each
$P_{k,t}$ is individually $I(1)$ because $F_t$ is a random walk. A stationary
linear combination of $I(1)$ series is cointegration by definition. \qed
\end{proof}

\begin{corollary}[Mean reversion and half-life]
The demeaned spread $\tilde Z_t:=Z_t-(a_k-a_\ell)$ obeys an AR(1) with
coefficient $\rho$, so it mean-reverts with half-life
\begin{equation}
t_{1/2}=\frac{\ln(1/2)}{\ln|\rho|}.
\end{equation}
\end{corollary}

\noindent
The package function \texttt{metrics.cointegration\_check} estimates $\rho$ by
OLS and reports $t_{1/2}$; on the default calibration ($\rho=0.92$) the
predicted half-life is $\ln(0.5)/\ln(0.92)\approx 8.3$ days, matching the
simulation output ($8.2$ days).

\subsection{Freight and basis}

Freight is volatile but does not trend without bound; we model log-freight as a
seasonal mean-reverting (Ornstein--Uhlenbeck/AR(1)) process:
\begin{equation}
\ln f_{ij,t}=\ln \bar f_{ij}+A\sin\!\Big(\tfrac{2\pi t}{252}\Big)+g_{ij,t},
\qquad g_{ij,t}=\phi\,g_{ij,t-1}+\xi_{ij,t}.
\end{equation}
The local \emph{basis} relative to the benchmark future is
$b_{k,t}:=P_{k,t}-H_t$, where $H_t=\bar H+F_t+c\,t$ embeds a deterministic carry
$c$. Because $H_t$ shares the world factor $F_t$, the basis $b_{k,t}=a_k+u_{k,t}
-\bar H-c\,t$ inherits stationarity around its (carry-adjusted) mean.

%======================================================================
\section{Distribution of the arbitrage margin}
\label{sec:dist}

We can now describe the margin's law. Hold FX at $e_{i,t}\approx 1$ and write
the transfer cost as
$K_{ij,t}=h^o_i+f_{ij,t}+\theta_{ij}+h^d_j+\text{(financing)}$.

\begin{proposition}[Margin decomposition]
\label{prop:margindist}
The arbitrage margin admits the decomposition
\begin{equation}
M_{ij,t}=\underbrace{(a_j-a_i)-\bar K_{ij}}_{\text{structural mean }\bar M_{ij}}
\;+\;\underbrace{(u_{j,t}-u_{i,t})}_{\text{stationary spread}}
\;-\;\underbrace{(f_{ij,t}-\bar f_{ij})}_{\text{freight shock}}
\;+\;(\text{financing, tariff, FX terms}),
\end{equation}
where $\bar K_{ij}=\E[K_{ij,t}]$. Consequently $M_{ij,t}$ is stationary; the
common world trend $F_t$ does \emph{not} enter, so margins do not inherit the
nonstationarity of price levels.
\end{proposition}

\begin{proof}
Substitute (\ref{eq:price}) and Definition~\ref{def:delivered} into
(\ref{eq:margin}); the $F_t$ terms in $P^d_{j,t}$ and $P^o_{i,t}$ cancel exactly
as in Proposition~\ref{prop:coint}. The remaining components are stationary
(AR(1) spreads, mean-reverting freight) or deterministic, so $M_{ij,t}$ is
stationary. \qed
\end{proof}

\begin{corollary}[Window-open probability]
If the stationary part of $M_{ij,t}$ is approximately Gaussian with mean
$\bar M_{ij}$ and standard deviation $\sigma_{M_{ij}}$, the long-run fraction of
days the window is open is
\begin{equation}
\Pr(M_{ij,t}>0)\;\approx\;\Phi\!\Big(\tfrac{\bar M_{ij}}{\sigma_{M_{ij}}}\Big),
\end{equation}
where $\Phi$ is the standard normal CDF. Structurally short, cheap lanes
($\bar M_{ij}>0$) are open most of the time; structurally unprofitable lanes
($\bar M_{ij}<0$) open only on favourable spread/freight shocks.
\end{corollary}

This corollary is the quantitative content of ``the arb is mostly closed'': for
a lane with $\bar M_{ij}<0$, windows are tail events driven by transitory
dislocations in the basis spread or freight.

%======================================================================
\section{Trading under transit risk and the futures hedge}
\label{sec:risk}

A cargo booked at $t$ is sold only on arrival at $t+\tau_{ij}$. The
\emph{realised} per-tonne PnL is therefore
\begin{equation}
\Pi_{ij,t}=P^d_{j,\,t+\tau_{ij}}-C_{ij,t},
\end{equation}
which differs from the signal $M_{ij,t}=P^d_{j,t}-C_{ij,t}$ by the destination
price move during transit, $P^d_{j,t+\tau}-P^d_{j,t}$. Since $P^d_j$ contains
the random walk $F$, this move has variance growing in $\tau$ --- the
\emph{flat-price} risk of a physical position.

\begin{proposition}[Futures hedge isolates the spatial spread]
\label{prop:hedge}
Sell one tonne of the benchmark future at $H_t$ when booking the cargo and buy
it back at arrival. The hedged PnL is
\begin{equation}
\Pi^{h}_{ij,t}=\bigl(P^d_{j,t+\tau}-C_{ij,t}\bigr)+\bigl(H_t-H_{t+\tau}\bigr)
= \underbrace{\bigl(b^d_{j,t+\tau}-(C_{ij,t}-H_t)\bigr)}_{\text{basis / spatial spread}}
\;-\;c\,\tau,
\end{equation}
so the world-factor term $F$ cancels between $P^d_j$ and $H$, leaving a
stationary basis spread plus a deterministic carry. The variance of $\Pi^h$ does
\emph{not} grow with $\tau$, unlike the unhedged $\Pi$.
\end{proposition}

\begin{proof}
Write $P^d_{j,t+\tau}=H_{t+\tau}+b^d_{j,t+\tau}$ by definition of basis.
Substituting into $\Pi^h$ cancels $H_{t+\tau}$:
$\Pi^h=H_{t+\tau}+b^d_{j,t+\tau}-C_{ij,t}+H_t-H_{t+\tau}
=b^d_{j,t+\tau}-(C_{ij,t}-H_t)$. With $H_t=\bar H+F_t+ct$ the residual carry is
$-c\tau$; all $F$ terms have cancelled. Basis is stationary
(Section~\ref{sec:dgp}), so $\Var(\Pi^h)$ is bounded in $\tau$. \qed
\end{proof}

Proposition~\ref{prop:hedge} is the formal statement of why physical arbitrage
desks hedge flat price with futures and trade the \emph{basis}: the hedge
removes the unbounded random-walk risk and leaves exactly the mean-reverting
spatial spread the trader has an informational edge on. The simulator
(\texttt{simulator.ArbitrageSimulator}) implements both the unhedged and hedged
PnL so the variance reduction can be measured directly.

%======================================================================
\section{From theory to the \texttt{geoarb} package}
\label{sec:map}

\begin{table}[h]
\centering
\small
\begin{tabular}{@{}lll@{}}
\toprule
\textbf{Object in this paper} & \textbf{Equation} & \textbf{Code} \\
\midrule
Network, voyage time & Def.~1 & \texttt{geography.py} \\
Delivered cost $C_{ij,t}$ & (\ref{eq:delivered}) & \texttt{arbitrage.delivered\_cost} \\
Arbitrage margin $M_{ij,t}$ & (\ref{eq:margin}) & \texttt{arbitrage.ArbitrageEngine} \\
No-arbitrage band & Prop.~1 & \texttt{ArbitrageEngine.window\_stats} \\
Spatial LP \& duals & (\ref{eq:lp}), Thm.~1 & \texttt{equilibrium.SpatialEquilibrium} \\
Common-trend prices & (\ref{eq:price}) & \texttt{dgp.SyntheticDGP} \\
Cointegration / half-life & Prop.~3, Cor.~1 & \texttt{metrics.cointegration\_check} \\
Margin distribution & Prop.~4, Cor.~2 & \texttt{ArbitrageEngine.window\_stats} \\
Hedged PnL & Prop.~5 & \texttt{simulator.ArbitrageSimulator} \\
\bottomrule
\end{tabular}
\caption{Correspondence between the formal objects and the implementation.}
\end{table}

%======================================================================
\section{Caveats and what the synthetic laboratory does \emph{not} prove}

The simulator routinely reports high Sharpe ratios. This is a feature of the
\emph{idealised} environment, not a claim about real markets. The synthetic
laboratory deliberately omits: execution slippage and market impact;
counterparty and demurrage risk; discrete, lumpy vessel availability;
information asymmetry and adverse selection; credit and margin calls on the
futures leg; and the fact that real basis is only partially predictable. Each
omission, once restored, compresses the Sharpe toward the thin, hard-won
margins that characterise the real physical-trading business. The value of the
laboratory is pedagogical and structural: it lets one see \emph{exactly} which
mechanism (cointegration, the no-arb band, the LP duals, the futures hedge)
produces each observed feature, because every mechanism is written down in
closed form.

%======================================================================
\section*{References (selected, for further reading)}
\addcontentsline{toc}{section}{References}
\small
\begin{enumerate}[label={[\arabic*]},leftmargin=*]
\item T.~Takayama and G.~G.~Judge, \emph{Spatial and Temporal Price and
Allocation Models}, North-Holland, 1971.
\item P.~A.~Samuelson, ``Spatial Price Equilibrium and Linear Programming,''
\emph{American Economic Review}, 1952.
\item R.~F.~Engle and C.~W.~J.~Granger, ``Co-integration and Error Correction,''
\emph{Econometrica}, 1987.
\item C.~Pirrong, \emph{The Economics of Commodity Trading Firms}, Trafigura
white paper, 2014.
\item H.~Working, ``The Theory of Price of Storage,'' \emph{American Economic
Review}, 1949.
\end{enumerate}

\vfill
\noindent\rule{\linewidth}{0.4pt}\\
{\footnotesize\textbf{Disclaimer.} All prices, freight rates, FX, and flows in
the companion software are synthetic and generated by the process of
Section~\ref{sec:dgp}. This document is educational and is not investment,
financial, or trading advice.}

\end{document}
